%\begin{abstract}
In questa esperienza di laboratorio si studiano i rivelatori a fotomoltiplicazione al Silicio (SiPM) e la loro applicazione nella rivelazione della radiazione $\beta^-$. \\ La prima parte consiste nella caratterizzazione di un SiPM a 667 celle, la quale comprende l'analisi del rumore intrinseco attraverso la misura del \textit{dark count rate} (DCR) e dell'\textit{optical cross talk} (OCT), nonché lo studio della loro dipendenza dalla temperatura. \\ Successivamente, mediante l’utilizzo di una sorgente luminosa regolabile (LED), si verifica la linearità dell'amplificazione del segnale, si stima il guadagno del SiPM verificandone la dipendenza dalla tensione di alimentazione e si studia lo spettro del LED in funzione della sua intensità.\\ Nella seconda parte dell’esperienza si analizza la rivelazione della radiazione $\beta^-$ usando uno scintillatore plastico accoppiato a un SiPM a 14400 celle; ponendo in particolare l'attenzione alla calibrazione in energia dell'apparato sperimentale e all'analisi degli spettri di sorgenti radioattive come $^{90}$Sr e $^{137}$Cs. \\ Con la sorgente di $^{90}$Sr, si studia inoltre la capacità di assorbimento della radiazione $\beta^-$ da parte di diversi materiali (carta e alluminio) stimandone caratteristiche fisiche quali il coefficiente di assorbimento e lo spessore. \\ Infine, viene utilizzato il sistema scintillatore+SiPM per la rivelazione dei muoni cosmici e per la verifica della legge di Bethe-Bloch.
%\end{abstract}
